\subsubsection{Statistical Models for Order Flow}

\begin{definition} \hlt{Point-Process View}\\
Order-book events can be modelled as point processes: random arrival times for insertions, cancellations, and transactions. The intensity $\lambda$ is the expected event rate over a short interval.
\end{definition}

\begin{taxonomy} \hlt{Order-Book Event Segmentation}
\begin{enumerate}[label=\roman*.]
\setlength{\itemsep}{0pt}
\item $\lambda^+$: limit-order insertions that add liquidity;
\item $\lambda^c$: cancellations that remove resting liquidity;
\item $\lambda^x$: transactions that consume resting liquidity.
\end{enumerate}
\end{taxonomy}

\begin{model} \hlt{Homogeneous Poisson Baseline}\\
The simplest model treats $\lambda^+$, $\lambda^c$, and $\lambda^x$ as constant and independent of the book state. It is a useful null model, but it cannot reproduce realistic queue-size distributions because event rates are strongly state-dependent.
\end{model}

\begin{mechanism} \hlt{Conditioning Improves Queue Dynamics}\\
Queue-reactive models condition event intensities on current and neighbouring queue sizes. Small best queues are easier to consume; large queues can trigger nearby cancellations.
\end{mechanism}

\begin{taxonomy} \hlt{Conditioning Segmentation}
\begin{enumerate}[label=\roman*.]
\setlength{\itemsep}{0pt}
\item \hlt{unconditional}: event type only;
\item \hlt{own queue}: event type and current queue size;
\item \hlt{neighbour queues}: own queue plus adjacent bid and ask queues;
\item \hlt{event time}: count events rather than seconds when clock time is less informative.
\end{enumerate}
\end{taxonomy}
